\section{Sharpness and three hyperbolic limit cycles}
\label{sec:sharpness}

To attain the upper bound, we need three displacement zeros whose
orbits approach the whole graphic. Two arise on a fixed compact
interval of normalized exit radii: along a buffered coefficient path,
the difference of the boundary actions gives a profile with a strict
negative minimum, and varying \(A\) sets its additive constant to a
small positive value. This produces two zeros. The third comes from
the organizing contraction at a larger central height, beyond the
compact exit interval. We must then show that this third height also
tends to zero and that all three orbits converge to the whole graphic.
These are the final steps of the proof below.

\subsection{A buffered path and its action difference}

Choose the coefficient path
\begin{equation}
\label{eq:sharpness-coefficient-path}
 \nu=\min\{1,(C_0-1)/64\},\qquad a=(1+\nu)/4,\qquad
 q_0=-as^5,\quad q_1=s^3,\quad q_2=-s,\qquad s>0.
\end{equation}
It lies in the buffered region \eqref{eq:direct-buffered-region}
with \(b=1\), \(r=s\), and \(b^2-4a=-\nu<0\).
The coefficient change \eqref{eq:direct-coefficient-bridge}
determines \(\mu(s)=(\mu_1(s),\mu_2(s),\mu_3(s))\to0\).
We keep \(C=C_0\) and write
\[
 p_s(A)=(A,C_0,\mu(s)).
\]
This path fixes the boundary drift through \(\mu\), while \(A\)
remains free to adjust the connection splitting. In particular, the
boundary actions are independent of this adjustment.

Besides the limiting exterior multiplier \(\Lambda_0\) in
\eqref{eq:exterior-multiplier-new}, put
\[
 \Lambda_{\rm b}=\exp(\pi/\sqrt\nu).
\]
The choice of \(\nu\) gives the multiplier margin
\begin{equation}
\label{eq:sharpness-multiplier-margin}
 \Lambda_0\Lambda_{\rm b}^2\ge \e^{3\pi/2}>100.
\end{equation}
Indeed, the logarithm of the left-hand side is
\(12\pi/\sqrt{C_0-1}\) when \(C_0-1\le64\), and
\(2\pi-4\pi/\sqrt{C_0-1}\) otherwise.

The following lemma identifies \(\Lambda_{\rm b}\) as the limit of the
boundary-action multiplier. In the comparison with
\cite[(2.13) and (6.8)]{HK2026} in
\cref{sec:direct-buffered-compact-primary}, their
\((\lambda_0^{\rm HK},\lambda_1^{\rm HK})=(-a,-1)\) give
\(\Lambda_{\rm b}=\e^{-K_{\rm HK}}\).

\begin{samepage}
\begin{lemma}[Buffered action limit]
\label{lem:buffered-action-limit}
Along \eqref{eq:sharpness-coefficient-path}, the constants
\(\mathfrak c_\sigma\) in \eqref{eq:direct-buffered-primary-action} satisfy
\[
 \mathfrak c_+(s)-\mathfrak c_-(s)\longrightarrow\frac{\pi}{\sqrt\nu},
 \qquad
 \Lambda_s:=\e^{\mathfrak c_+(s)-\mathfrak c_-(s)}\longrightarrow\Lambda_{\rm b}.
\]
They are independent of \(A\).
\end{lemma}
\end{samepage}

\begin{proof}
Here \(D_\sigma\), \(\mathcal A_\sigma\), \(L_\sigma\), \(\mathfrak c_\sigma\) and the
reference radius \(v_b\) are those of \cref{sec:direct-buffered-negative}
with \(r=s\) and \(b=1\); thus
\(D_\sigma(t)=as^2-\sigma st+t^2-\sigma t^3R_q(\sigma st)\).
The fixed positive margin \(\nu\) gives
\(D_\sigma(t)\ge c_\nu(s^2+t^2)\) for \(0\le t\le v_b\).
The same bound holds after replacing \(R_q\) by \(-1\).
Since \(R_q(\sigma st)=-1+O(s)\), this replacement changes
each action by at most
\[
 Cs\int_0^{v_b}\frac{t^4}{(s^2+t^2)^2}\dd t=O(s).
\]

We separate the remaining integral at a fixed
\(0<\varepsilon<v_b\) to distinguish the action accumulated near
the central section from the outer reference terms. On
\([0,\varepsilon]\), deleting the
remaining cubic term changes each action by at most \(C\varepsilon\).
With \(t=sx\), the difference of the resulting inner actions tends to
\begin{align}
\label{eq:sharpness-action-integral}
 c_\infty
 &=\int_0^\infty x
   \left(\frac1{x^2-x+a}-\frac1{x^2+x+a}\right)\dd x \notag\\
 &=\int_0^\infty
       \frac{2x^2}{x^4+(2a-1)x^2+a^2}\dd x
   =\frac{\pi}{\sqrt{4a-1}}.
\end{align}
For the last equality, substitute \(x=\sqrt a\,u\), use
\(u\mapsto1/u\) to equate the integrals with numerators \(u^2\)
and \(1\), and apply \(y=u-1/u\) to their sum.
Both ends of the integral converge.

On \([\varepsilon,v_b]\), the action difference converges to
\[
 (L_+-L_-)(v_b)-(L_+-L_-)(\varepsilon).
\]
Subtract the reference terms in \(\mathfrak c_+-\mathfrak c_-\),
first let \(s\to0\), and then let \(\varepsilon\to0\).
The remaining outer term vanishes because
\((L_+-L_-)(\varepsilon)=\log((1-\varepsilon)/(1+\varepsilon))\to0\).
This proves the result.
\end{proof}

\subsection{Half-orbits with a prescribed exit}

The action comparison must be realized by half-orbits with a common
central height. We first construct each half-orbit backward from its
exit; monotonicity of the resulting height then gives the matching.
Fix \(J_+=[\rho_{\rm lo},\rho_{\rm hi}]\Subset(0,\infty)\) and
\(J_-\Subset(0,1)\) whose interior contains
\(f_{\Lambda_{\rm b}}(J_+)\). For a terminal radius
\(x_e\in J_\sigma\), prescribe
\(\eta=\sigma sx_e\), \(z=s^3\delta_{\rm an}x_e^2\) and integrate
toward the central section. In the local variables \(t=|\eta|/s\)
and \(y=z/s^3\), the exact radial equation is
\[
 y_t=\frac{ty}{s\{D_\sigma(t)+B(\sigma st,s^3y,p_s(A))y\}},
 \qquad y(x_e)=\delta_{\rm an}x_e^2.
\]
The quadratic margin and the upper end of \(J_-\) give
\(c(s^2+t^2)\le D_\sigma(t)\le C(s^2+t^2)\) on these fixed
intervals. The backward solution remains positive and bounded,
and reaches \(t=0\) with \(h_\sigma(x_e)=s^3y(0)>0\).
Its radius and height remain in the normal-rate box for small \(s\),
and \(q-zB<0\) identifies it with the corresponding half-orbit
of the original field.

The correction to its boundary action is exactly
\[
 \mathscr E_\sigma(x_e)
 =\int_0^{x_e}\frac{tBy}{sD_\sigma(D_\sigma+By)}\,dt
 =\int_{y(0)}^{y(x_e)}\frac{B}{D_\sigma}\,dy\ge0,
\]
with the coefficients evaluated along the half-orbit.
Choose a fixed \(\kappa>0\) with \(2\kappa<\inf(J_+\cup J_-)\).
The logarithmic growth rate on \([\kappa,2\kappa]\) is at least
\(c/s\), so \(y(\kappa)\le C\e^{-c/s}\).
The action correction up to \(\kappa\) is at most
\(Cs^{-2}y(\kappa)\); thereafter it is at most \(Cy(x_e)\),
since \(D_\sigma\ge c\kappa^2\). Thus
\(\mathscr E_\sigma\) is uniformly bounded and
\begin{equation}
\label{eq:sharpness-common-source}
 s\log(h_\sigma(x_e)/s^3)=-\mathcal A_\sigma(x_e)
                   +sM_\sigma^0(x_e),\qquad
 M_\sigma^0=\log(\delta_{\rm an}x_e^2)+\mathscr E_\sigma,
 \qquad \|M_\sigma^0\|_{C^0}\le C.
\end{equation}
All constants may depend on \(J_\pm,C_0\) and the fixed \(A\)-interval;
\(\delta_{\rm an}\) is held fixed before reducing \(s\).

At \(t=K_0s\) the normalized height
\(y/t^2\) is exponentially small. At any later crossing of
\(y/t^2=\delta_{\rm an}\),
\[
 \frac{d\log(y/t^2)}{d\log t}
 =\frac{t^2}{s(D_\sigma+B\delta_{\rm an}t^2)}-2>0
\]
for small \(s\), so the prescribed point is the first upward crossing
after \(t=K_0s\).
The central height decreases strictly with its exit radius:
the radial solution has positive derivative in its initial height,
whereas its terminal radial derivative exceeds the derivative of
\(\delta_{\rm an}x_e^2\) by the same inequality.

On \(J_\pm\), subtraction of the action integrands from their limits
gives \(\mathcal A_\sigma=\mathfrak c_\sigma+L_\sigma+O(s)\).
The strict endpoint margins of \(J_-\),
\cref{lem:buffered-action-limit}, and
\eqref{eq:sharpness-common-source} therefore bracket
\(h_+(\rho)\) between the two negative endpoint heights.
Continuity and strict decrease give a unique negative exit \(\omega\)
with the same central height. Equating their actions and using the
positive lower bound on \(L_-'\) gives
\begin{equation}
\label{eq:sharpness-negative-exit}
 \omega=f_{\Lambda_s}(\rho)+O_{C^0}(s),\qquad
 f_\Lambda(\rho)=\frac{\Lambda\rho}{1+(1+\Lambda)\rho}.
\end{equation}
Uniqueness at the regular central point makes these two half-orbits
parts of one orbit. The continuation below places their endpoints
in the fixed windows.

\subsection{A transverse organizing splitting}

Write \(p_0(A)=(A,C_0,0,0,0)\). The tuning argument requires a
strict sign for the derivative of the organizing splitting. We first
construct the two boundary graphs and their dependence on \(A\).
The variational equation along the parabola gives the sign of their
splitting in an affine coordinate; the section change has positive
derivative and preserves this sign in the endpoint coordinates.

\begin{lemma}[Transversality of the organizing splitting]
\label{lem:organizing-splitting-transversality}
For \(A\) close to \(1\), the graph equation
\eqref{eq:physical-graph-equation-new} at \(p_0(A)\) has unique
solutions
\[
 U_\sigma^0(z;A)=z+O(z^{3/2}),\qquad \sigma=\pm,
\]
which depend \(C^1\) on \(A\). Set
\[
 b_\sigma^0(A)=z_*-U_\sigma^0(z_*;A),\qquad
 \mathcal S_0(A)=\Pi_{p_0(A)}(b_-^0(A))-b_+^0(A).
\]
Then \(\mathcal S_0(1)=0\) and \(\mathcal S_0'(1)<0\).
\end{lemma}

\begin{proof}
\emph{Boundary graphs and their variations.}
At \(\mu=0\), both \(q\) and \(\eta=v+v^2+C_0vz\) are independent of \(A\).
Differentiating the exact projective field gives
\begin{equation}
\label{eq:sharpness-B-variation}
 B_A=-v(1+2v+C_0z),\qquad B(0,z,p_0(A))=1+C_0z.
\end{equation}
On \(U\asymp z\), uniformly for \(A\) near \(1\), the graph
right side \(\mathcal G_{\sigma,p_0(A)}\) therefore satisfies
\[
 \mathcal G_{\sigma,p_0(A)}=1+O(\sqrt z),\qquad
 |\mathcal G_{\sigma,U}|\le Cz^{-1/2},\qquad
 |\mathcal G_{\sigma,A}|\le C\sqrt z,
\]
where the middle bound is \eqref{eq:direct-tail-coefficients} at
\(\mu=0\).
The Volterra operator
\[
 \mathcal T_AU(z)=\int_0^z\mathcal G_{\sigma,p_0(A)}(U(t),t)\dd t
\]
maps \(|U(z)/z-1|\le1/2\) into itself. In the norm
\(\|U\|_*=\sup_{0<z\le z_*}|U(z)|/z\), its Lipschitz
constant is at most \(C\sqrt{z_*}<1\).
The operator is \(C^1\) in \(U,A\) on a slightly wider positive
cone, so differentiation of the convergent iteration gives
the asserted \(C^1\) fixed point. Its parameter derivative satisfies
\begin{equation}
\label{eq:sharpness-origin-variation}
 \begin{gathered}
 (\partial_AU_\sigma^0)'=\mathcal G_U\partial_AU_\sigma^0+B_A,\qquad
 \partial_AU_\sigma^0(0)=0,\\
 \partial_AU_\sigma^0(z)=O(z^{3/2}).
 \end{gathered}
\end{equation}

\emph{Comparison along the organizing parabola.}
For the entire \(p_0(A)\)-family, the function \(H_0\) of
\eqref{eq:parabola-invariance-new} satisfies the exact identity
\begin{equation}
\label{eq:sharpness-affine-variation}
 \dot H_0=2xH_0-(A-1)x^2.
\end{equation}
Continue the left boundary graph forward from the negative section,
and the right boundary graph backward from the positive section,
to the same point \(x\) on the organizing parabola. Denote their
\(A\)-derivatives at \(A=1\), measured in the \(H_0\)-coordinate
at fixed \(x\), by \(\mathfrak k_{\rm L}\) and
\(\mathfrak k_{\rm R}\), respectively. Their difference is the
signed variation of the splitting at that point. Changing from
fixed-\(x\) to fixed-\(z\) sections adds a tangential variation
along the organizing parabola, whose contribution to the variation
of \(H_0\) at \(A=1\) is zero.

Use \(x\) as time, as in the proof of
\cref{lem:exterior-multiplier-new}. The last term of
\eqref{eq:sharpness-affine-variation} gives the forcing in the
variational equation. With \(P_{\rm aff}(x)=x^2+2x+C_0>0\), both
variations satisfy
\[
 \mathfrak k'=\frac{4x}{P_{\rm aff}(x)}\mathfrak k-\frac{2x^2}{P_{\rm aff}(x)}.
\]
The boundary conditions follow from
\eqref{eq:sharpness-origin-variation}: at fixed \(z\),
\[
 \partial_Av
   =\frac{\sigma\partial_AU_\sigma^0}
       {\sqrt{2U_\sigma^0}(1+2v+C_0z)}=O(z),\qquad
 \partial_AH_0=-\frac{v\,\partial_Av}{z^2}
       =O(z^{-1/2})=O(|x|).
\]
With
\[
 W(x)=\exp\left(-\int_0^x\frac{4t}{P_{\rm aff}(t)}\dd t\right)>0,
\]
direct integration gives \(W(x)\asymp|x|^{-4}\) at both ends.
Hence \(W\mathfrak k_{\rm L}\to0\) as \(x\to-\infty\) and
\(W\mathfrak k_{\rm R}\to0\) as \(x\to+\infty\). Variation of constants
with these two endpoint conditions yields
\begin{equation}
\label{eq:sharpness-positive-integral}
 \mathfrak k_{\rm L}(x)-\mathfrak k_{\rm R}(x)
 =-W(x)^{-1}\int_{-\infty}^{\infty}
                  \frac{2W(t)t^2}{P_{\rm aff}(t)}\dd t<0.
\end{equation}
The integral is absolutely convergent and strictly positive.

\emph{The splitting in endpoint coordinates.}
To relate this variation to the displacement, let \(\chi_\pm\)
be the \(H_0\)-coordinates on the two fixed-\(z\) sections.
The section changes are regular, so \(\chi_\pm'(I_\pm^0)>0\), and
the chain rule at the positive section gives
\begin{equation}
\label{eq:sharpness-splitting-bridge}
 \mathcal S_0'(1)
 =\frac{1}{\chi_+'(I_+^0)}
   \bigl(\mathfrak k_{\rm L}(x_+)-\mathfrak k_{\rm R}(x_+)\bigr),
 \qquad \chi_+'(I_+^0)>0 .
\end{equation}
Here \(x_+\) is the organizing point on the positive section, where
both variations are evaluated. The finite exterior propagation is
included in \(\mathfrak k_{\rm L}(x_+)\), and
\(1/\chi_+'(I_+^0)\) is the positive \(H_0\)-to-\(I\) endpoint
conversion factor computed in \cref{lem:exterior-multiplier-new}.
Thus \eqref{eq:sharpness-positive-integral} and
\eqref{eq:sharpness-splitting-bridge} give \(\mathcal S_0'(1)<0\).
\end{proof}

\subsection{Tuning the splitting and obtaining three cycles}

\begin{proof}[Proof of \cref{prop:three-cycle-lower-bound}]
\emph{Endpoint preparation.}
We first continue the half-orbits from the compact exit intervals to
the fixed sections, where their endpoint maps can be composed with
the exterior transition. Use
\cref{lem:direct-small-height-continuation} with
\(\ell=\delta=s\) and \(\zeta_0(\rho_\sigma)=\delta_{\rm an}\rho_\sigma^2\),
as in \eqref{eq:direct-buffered-primary-endpoints}. On the path
\eqref{eq:sharpness-coefficient-path}, the coefficient ratios are
bounded, and the fixed compact exit intervals have positive drift
margins, uniformly for \(A\) near \(1\).
Choose the large height \(Z_1=Z_0\) of that lemma uniformly in \(A\),
and put \(t_s=s^2Z_0\). The continued and reference energies lie in
the common cone. Let \(U_\sigma^{\rm ref}\) be the tail with initial
energy \(U_\sigma^{\rm ref}(t_s)=B_0(p)t_s\), and define
\begin{equation}
\label{eq:sharpness-continuous-tail}
 \begin{aligned}
 b_\sigma^e(p)&=z_*-U_\sigma^{\rm ref}(z_*),\\
 a_\sigma^e(p)&=
 \left.\partial_uU_\sigma(z_*;t_s,u,p)\right|_{u=B_0(p)t_s}>0.
 \end{aligned}
\end{equation}
The reference endpoint and the tail-energy derivative are smooth in \(A\)
for fixed \(s>0\). They are the lemma's \(b_\sigma\) and
\(\varpi_{t_s}A_\sigma\), respectively.
Since \(\varpi_{t_s}=1+O(s(1+|\log s|))\), the coefficients
\(a_\sigma^e\) have uniform positive bounds, and
\eqref{eq:direct-small-height-endpoint} gives
\begin{equation}
\label{eq:sharpness-prepared-endpoints}
 I_\sigma(\rho_\sigma;p)
 =b_\sigma^e-\tfrac12s^2a_\sigma^e\rho_\sigma^2
                +s^2O_{C^3}(s(1+|\log s|)).
\end{equation}

\emph{Domains of the exterior transition.}
To place the reference and continued endpoints in the exterior windows
for small \(s\), uniformly on a fixed small \(A\)-interval, compare
their tails to \(U_\sigma^0(z;A)\). On the common cone, the difference
of the graph right sides at \(p_s(A)\) and \(p_0(A)\) is bounded by
\[
 C(s+s^5/z+s^3/\sqrt z),
\]
whereas, by \eqref{eq:direct-tail-coefficients}, a Lipschitz
coefficient for the difference is
\begin{equation}
\label{eq:sharpness-tail-comparison}
 2s/z+C(z^{-1/2}+s^3z^{-3/2}).
\end{equation}
Its integral over \([t_s,z_*]\) is uniformly bounded.
The initial energy difference is \(O(s^2)\), and the integrated
forcing is \(O(s)\). Gronwall's inequality gives, uniformly in
\(A\) and the fixed exit ranges,
\begin{equation}
\label{eq:sharpness-endpoint-limit}
 b_\sigma^e(A,s)\longrightarrow b_\sigma^0(A),\qquad
 I_\sigma(\rho_\sigma;A,s)\longrightarrow b_\sigma^0(A).
\end{equation}
First restrict \(A\) so that \(b_\sigma^0(A)\in\operatorname{int}W_\sigma\).
Uniform convergence then puts, for small \(s\), all the endpoints in
\eqref{eq:sharpness-endpoint-limit}, including those at the wider
comparison exits, in \(\operatorname{int}W_\sigma\). The continued
half-orbits thus give admissible central heights in
\(\mathcal H_p\). All these endpoints also lie in the enlarged
windows, where the exterior map and its derivatives are defined for
the displacement expansion and splitting comparison below.

\emph{The normalized displacement.}
We separate the constant term, which will be adjusted by \(A\), from
the coefficient of the compact exit profile. Define the continuous
splitting and slope coefficient
\begin{equation}
\label{eq:sharpness-actual-splitting}
 \mathcal S_s(A)=\Pi_{p_s(A)}(b_-^e)-b_+^e,\qquad
 \lambda_s(A)=\Pi_{p_s(A)}'(b_-^e)\frac{a_-^e}{a_+^e}.
\end{equation}
By \eqref{eq:sharpness-negative-exit} and
\eqref{eq:sharpness-prepared-endpoints}, the prepared displacement
\eqref{eq:direct-buffered-prepared-displacement} holds with \(r=s\),
\(\Lambda=\Lambda_s\), \(\lambda(p)=\lambda_s(A)\) and
\(D(p)=2\mathcal S_s(A)/(s^2a_+^e)\):
\begin{equation}
\label{eq:sharpness-prepared-displacement}
 \frac{2\Delta_{p_s(A)}(h(\rho;p_s(A)))}{s^2a_+^e}
 =\frac{2\mathcal S_s(A)}{s^2a_+^e}
      +g_{\lambda_s(A),\Lambda_s}(\rho)
      +O_{C^0}(s(1+|\log s|)),
\end{equation}
where
\[
 g_{\lambda,\Lambda}(\rho)
   =\rho^2-\lambda f_\Lambda(\rho)^2
   =\rho^2-\frac{\lambda\Lambda^2\rho^2}
                         {[1+(1+\Lambda)\rho]^2}.
\]

The three test radii must work throughout the tuning interval in
\(A\). We therefore bound the slope coefficient uniformly first.
The common factor
\(\varpi_{t_s}\) cancels from the quotient, so
\eqref{eq:direct-tail-first-slope} gives
\(\log(a_-^e/a_+^e)=O(\sqrt{z_*})+o_s(1)\). Also,
\(\Pi_0'(I_-^0)=\Lambda_0\{1+O(\sqrt{z_*})\}\) by
\cref{lem:exterior-multiplier-new}. As in
\cref{sec:direct-passage-estimates}, the coefficient of \(\sqrt{z_*}\)
depends only on the common cone; constants depending on the fixed
exit ranges enter only the \(o_s(1)\) term. Thus we may first choose
the fixed section, then restrict \(A\) and \(s\), so that
\begin{equation}
\label{eq:sharpness-lambda-range}
 \lambda_{\rm lo}:=\Lambda_0/2
 \le\lambda_s(A)\le
 \lambda_{\rm hi}:=(1+\Lambda_0)/2<1.
\end{equation}

\emph{Three test radii.}
The multiplier margin makes the model negative at an interior
minimum; a small positive additive constant will give the two outer
positive signs. To choose radii with fixed sign margins, put
\[
 M_{\rm b}=\sqrt{\lambda_{\rm lo}}\Lambda_{\rm b}>1,\qquad
 K_{\rm b}=1+\Lambda_{\rm b}.
\]
The strict negative minimum of
\(g_{\lambda_{\rm lo},\Lambda_{\rm b}}\) occurs at
\begin{equation}
\label{eq:sharpness-negative-minimum}
 \rho_{\rm b}=\frac{M_{\rm b}^{2/3}-1}{K_{\rm b}},\qquad
 g_{\lambda_{\rm lo},\Lambda_{\rm b}}(\rho_{\rm b})
       =-\frac{(M_{\rm b}^{2/3}-1)^3}{K_{\rm b}^2}<0.
\end{equation}
This follows directly from
\[
 g_{\lambda,\Lambda}'(\rho)
   =2\rho\left(1-\frac{\lambda\Lambda^2}
                            {[1+(1+\Lambda)\rho]^3}\right).
\]
Choose \(0<D_*<-g_{\lambda_{\rm lo},\Lambda_{\rm b}}(\rho_{\rm b})/4\).
Since \(g\) decreases in \(\lambda\),
\(D_*+g_{\lambda,\Lambda_{\rm b}}(\rho_{\rm b})<-3D_*\)
throughout \eqref{eq:sharpness-lambda-range}.
Take \(0<\rho_{\rm a}<\rho_{\rm b}\) so small that
\[
 \sup_{\lambda_{\rm lo}\le\lambda\le\lambda_{\rm hi}}
       |g_{\lambda,\Lambda_{\rm b}}(\rho_{\rm a})|<D_*/4,
\]
and then take \(\rho_{\rm c}>\rho_{\rm b}\) with
\(\rho_{\rm c}^2>\lambda_{\rm hi}+1\).
Since \(0<f_\Lambda<1\), the last choice gives
\(g_{\lambda,\Lambda}(\rho_{\rm c})>1\).
By \cref{lem:buffered-action-limit}, these choices imply,
uniformly in \(A\), for small \(s\),
\begin{equation}
\label{eq:sharpness-three-model-signs}
 \begin{aligned}
 D_*+g_{\lambda_s(A),\Lambda_s}(\rho_{\rm a})&>D_*/2,\\
 D_*+g_{\lambda_s(A),\Lambda_s}(\rho_{\rm b})&<-2D_*,\\
 D_*+g_{\lambda_s(A),\Lambda_s}(\rho_{\rm c})&>1+D_*.
 \end{aligned}
\end{equation}
All three radii and their negative comparison intervals are
fixed before the final reduction of \(s\). Apply the compact construction
above with \(J_+\) containing \([\rho_{\rm a},\rho_{\rm c}]\) and
\(f_{\Lambda_{\rm b}}(J_+)\subset\operatorname{int}J_-\).
\Cref{fig:sharpness-model}
shows the model profile and the three radii for \(C_0=65\).

\begin{figure}[t]
\centering
% Generated by scripts/figure_sharpness_model.py; do not edit.
% C_0 = 65, lambda = Lambda_0 = 0.207880, Lambda = Lambda_b = 23.140693,
% rho_a = 0.00251189, rho_b = 0.116736, rho_c = 1.3, D_* = 0.0076805.
\newcommand{\figModelFull}[1]{\draw[#1] (0.000,0.013) -- (0.014,0.011) -- (0.027,0.007) -- (0.041,0.001) -- (0.054,-0.006) -- (0.068,-0.014) -- (0.081,-0.021) -- (0.095,-0.029) -- (0.108,-0.036) -- (0.121,-0.043) -- (0.135,-0.050) -- (0.149,-0.057) -- (0.162,-0.064) -- (0.176,-0.070) -- (0.189,-0.075) -- (0.203,-0.081) -- (0.216,-0.086) -- (0.230,-0.091) -- (0.243,-0.095) -- (0.257,-0.100) -- (0.270,-0.104) -- (0.284,-0.108) -- (0.297,-0.111) -- (0.310,-0.115) -- (0.324,-0.118) -- (0.338,-0.121) -- (0.351,-0.124) -- (0.365,-0.126) -- (0.378,-0.129) -- (0.392,-0.131) -- (0.405,-0.133) -- (0.419,-0.135) -- (0.432,-0.137) -- (0.446,-0.138) -- (0.459,-0.140) -- (0.472,-0.141) -- (0.486,-0.142) -- (0.500,-0.144) -- (0.513,-0.145) -- (0.527,-0.145) -- (0.540,-0.146) -- (0.553,-0.147) -- (0.567,-0.147) -- (0.581,-0.148) -- (0.594,-0.148) -- (0.608,-0.148) -- (0.621,-0.148) -- (0.635,-0.148) -- (0.648,-0.148) -- (0.662,-0.148) -- (0.675,-0.148) -- (0.689,-0.148) -- (0.702,-0.147) -- (0.716,-0.147) -- (0.729,-0.146) -- (0.743,-0.145) -- (0.756,-0.145) -- (0.770,-0.144) -- (0.783,-0.143) -- (0.797,-0.142) -- (0.810,-0.141) -- (0.824,-0.140) -- (0.837,-0.139) -- (0.851,-0.138) -- (0.864,-0.136) -- (0.877,-0.135) -- (0.891,-0.133) -- (0.905,-0.132) -- (0.918,-0.130) -- (0.932,-0.129) -- (0.945,-0.127) -- (0.959,-0.125) -- (0.972,-0.124) -- (0.986,-0.122) -- (0.999,-0.120) -- (1.012,-0.118) -- (1.026,-0.116) -- (1.040,-0.114) -- (1.053,-0.112) -- (1.067,-0.109) -- (1.080,-0.107) -- (1.094,-0.105) -- (1.107,-0.102) -- (1.121,-0.100) -- (1.134,-0.098) -- (1.148,-0.095) -- (1.161,-0.092) -- (1.175,-0.090) -- (1.188,-0.087) -- (1.202,-0.084) -- (1.215,-0.082) -- (1.229,-0.079) -- (1.242,-0.076) -- (1.256,-0.073) -- (1.269,-0.070) -- (1.282,-0.067) -- (1.296,-0.064) -- (1.310,-0.061) -- (1.323,-0.057) -- (1.337,-0.054) -- (1.350,-0.051) -- (1.364,-0.048) -- (1.377,-0.044) -- (1.391,-0.041) -- (1.404,-0.037) -- (1.417,-0.034) -- (1.431,-0.030) -- (1.445,-0.027) -- (1.458,-0.023) -- (1.472,-0.019) -- (1.485,-0.015) -- (1.499,-0.012) -- (1.512,-0.008) -- (1.526,-0.004) -- (1.539,0.000) -- (1.552,0.004) -- (1.566,0.008) -- (1.580,0.012) -- (1.593,0.016) -- (1.607,0.020) -- (1.620,0.025) -- (1.634,0.029) -- (1.647,0.033) -- (1.661,0.038) -- (1.674,0.042) -- (1.688,0.046) -- (1.701,0.051) -- (1.715,0.055) -- (1.728,0.060) -- (1.742,0.065) -- (1.755,0.069) -- (1.769,0.074) -- (1.782,0.079) -- (1.796,0.083) -- (1.809,0.088) -- (1.823,0.093) -- (1.836,0.098) -- (1.850,0.103) -- (1.863,0.108) -- (1.877,0.113) -- (1.890,0.118) -- (1.904,0.123) -- (1.917,0.128) -- (1.931,0.134) -- (1.944,0.139) -- (1.958,0.144) -- (1.971,0.150) -- (1.985,0.155) -- (1.998,0.160) -- (2.011,0.166) -- (2.025,0.171) -- (2.039,0.177) -- (2.052,0.182) -- (2.066,0.188) -- (2.079,0.194) -- (2.092,0.199) -- (2.106,0.205) -- (2.120,0.211) -- (2.133,0.217) -- (2.147,0.223) -- (2.160,0.229) -- (2.174,0.235) -- (2.187,0.241) -- (2.200,0.247) -- (2.214,0.253) -- (2.228,0.259) -- (2.241,0.265) -- (2.255,0.271) -- (2.268,0.278) -- (2.281,0.284) -- (2.295,0.290) -- (2.309,0.297) -- (2.322,0.303) -- (2.336,0.309) -- (2.349,0.316) -- (2.363,0.323) -- (2.376,0.329) -- (2.389,0.336) -- (2.403,0.342) -- (2.417,0.349) -- (2.430,0.356) -- (2.444,0.363) -- (2.457,0.369) -- (2.470,0.376) -- (2.484,0.383) -- (2.498,0.390) -- (2.511,0.397) -- (2.525,0.404) -- (2.538,0.411) -- (2.551,0.418) -- (2.565,0.426) -- (2.579,0.433) -- (2.592,0.440) -- (2.606,0.447) -- (2.619,0.455) -- (2.632,0.462) -- (2.646,0.469) -- (2.660,0.477) -- (2.673,0.484) -- (2.687,0.492) -- (2.700,0.499) -- (2.714,0.507) -- (2.727,0.515) -- (2.740,0.522) -- (2.754,0.530) -- (2.768,0.538) -- (2.781,0.546) -- (2.795,0.553) -- (2.808,0.561) -- (2.822,0.569) -- (2.835,0.577) -- (2.849,0.585) -- (2.862,0.593) -- (2.876,0.601) -- (2.889,0.609) -- (2.902,0.618) -- (2.916,0.626) -- (2.930,0.634) -- (2.943,0.642) -- (2.957,0.651) -- (2.970,0.659) -- (2.984,0.667) -- (2.997,0.676) -- (3.010,0.684) -- (3.024,0.693) -- (3.038,0.701) -- (3.051,0.710) -- (3.065,0.719) -- (3.078,0.727) -- (3.092,0.736) -- (3.105,0.745) -- (3.119,0.754) -- (3.132,0.762) -- (3.146,0.771) -- (3.159,0.780) -- (3.172,0.789) -- (3.186,0.798) -- (3.200,0.807) -- (3.213,0.816) -- (3.227,0.825) -- (3.240,0.835) -- (3.254,0.844) -- (3.267,0.853) -- (3.280,0.862) -- (3.294,0.872) -- (3.308,0.881) -- (3.321,0.890) -- (3.335,0.900) -- (3.348,0.909) -- (3.362,0.919) -- (3.375,0.928) -- (3.389,0.938) -- (3.402,0.947) -- (3.416,0.957) -- (3.429,0.967) -- (3.442,0.977) -- (3.456,0.986) -- (3.470,0.996) -- (3.483,1.006) -- (3.497,1.016) -- (3.510,1.026) -- (3.524,1.036) -- (3.537,1.046) -- (3.550,1.056) -- (3.564,1.066) -- (3.578,1.076) -- (3.591,1.086) -- (3.605,1.097) -- (3.618,1.107) -- (3.632,1.117) -- (3.645,1.128) -- (3.659,1.138) -- (3.672,1.148) -- (3.686,1.159) -- (3.699,1.169) -- (3.712,1.180) -- (3.726,1.190) -- (3.740,1.201) -- (3.753,1.212) -- (3.767,1.222) -- (3.780,1.233) -- (3.794,1.244) -- (3.807,1.255) -- (3.821,1.266) -- (3.834,1.277) -- (3.848,1.287) -- (3.861,1.298) -- (3.875,1.309) -- (3.888,1.320) -- (3.902,1.332) -- (3.915,1.343) -- (3.929,1.354) -- (3.942,1.365) -- (3.956,1.376) -- (3.969,1.388) -- (3.982,1.399) -- (3.996,1.410) -- (4.010,1.422) -- (4.023,1.433) -- (4.037,1.445) -- (4.050,1.456) -- (4.064,1.468) -- (4.077,1.479) -- (4.091,1.491) -- (4.104,1.503) -- (4.117,1.514) -- (4.131,1.526) -- (4.145,1.538) -- (4.158,1.550) -- (4.171,1.562) -- (4.185,1.573) -- (4.199,1.585) -- (4.212,1.597) -- (4.226,1.609) -- (4.239,1.622) -- (4.253,1.634) -- (4.266,1.646) -- (4.280,1.658) -- (4.293,1.670) -- (4.307,1.682) -- (4.320,1.695) -- (4.333,1.707) -- (4.347,1.719) -- (4.361,1.732) -- (4.374,1.744) -- (4.388,1.757) -- (4.401,1.769) -- (4.415,1.782) -- (4.428,1.795) -- (4.442,1.807) -- (4.455,1.820) -- (4.469,1.833) -- (4.482,1.845) -- (4.495,1.858) -- (4.509,1.871) -- (4.523,1.884) -- (4.536,1.897) -- (4.550,1.910) -- (4.563,1.923) -- (4.577,1.936) -- (4.590,1.949) -- (4.604,1.962) -- (4.617,1.975) -- (4.631,1.989) -- (4.644,2.002) -- (4.658,2.015) -- (4.671,2.028) -- (4.685,2.042) -- (4.698,2.055) -- (4.712,2.069) -- (4.725,2.082) -- (4.739,2.096) -- (4.752,2.109) -- (4.766,2.123) -- (4.779,2.136) -- (4.793,2.150) -- (4.806,2.164) -- (4.820,2.177) -- (4.833,2.191) -- (4.847,2.205) -- (4.860,2.219) -- (4.873,2.233) -- (4.887,2.247) -- (4.901,2.261) -- (4.914,2.275) -- (4.928,2.289) -- (4.941,2.303) -- (4.955,2.317) -- (4.968,2.331) -- (4.982,2.346) -- (4.995,2.360) -- (5.009,2.374) -- (5.022,2.388) -- (5.035,2.403) -- (5.049,2.417) -- (5.063,2.432) -- (5.076,2.446) -- (5.090,2.461) -- (5.103,2.475) -- (5.117,2.490) -- (5.130,2.505) -- (5.144,2.519) -- (5.157,2.534) -- (5.171,2.549) -- (5.184,2.563) -- (5.197,2.578) -- (5.211,2.593) -- (5.225,2.608) -- (5.238,2.623) -- (5.252,2.638) -- (5.265,2.653) -- (5.279,2.668) -- (5.292,2.683) -- (5.306,2.698) -- (5.319,2.714) -- (5.332,2.729) -- (5.346,2.744) -- (5.360,2.759) -- (5.373,2.775) -- (5.387,2.790) -- (5.400,2.806);}
\newcommand{\figModelZoom}[1]{\draw[#1] (7.200,0.230) -- (7.214,0.230) -- (7.228,0.230) -- (7.242,0.230) -- (7.256,0.230) -- (7.270,0.230) -- (7.284,0.230) -- (7.298,0.230) -- (7.312,0.230) -- (7.326,0.230) -- (7.340,0.230) -- (7.354,0.230) -- (7.368,0.230) -- (7.382,0.230) -- (7.396,0.230) -- (7.410,0.230) -- (7.424,0.230) -- (7.438,0.230) -- (7.452,0.230) -- (7.466,0.230) -- (7.480,0.230) -- (7.494,0.230) -- (7.508,0.230) -- (7.522,0.230) -- (7.536,0.230) -- (7.550,0.230) -- (7.564,0.230) -- (7.578,0.230) -- (7.592,0.230) -- (7.606,0.229) -- (7.620,0.229) -- (7.634,0.229) -- (7.648,0.229) -- (7.662,0.229) -- (7.676,0.229) -- (7.690,0.229) -- (7.704,0.229) -- (7.718,0.229) -- (7.732,0.229) -- (7.746,0.229) -- (7.760,0.229) -- (7.774,0.229) -- (7.788,0.229) -- (7.802,0.229) -- (7.816,0.229) -- (7.830,0.229) -- (7.844,0.229) -- (7.858,0.229) -- (7.872,0.229) -- (7.886,0.228) -- (7.900,0.228) -- (7.914,0.228) -- (7.928,0.228) -- (7.942,0.228) -- (7.956,0.228) -- (7.970,0.228) -- (7.984,0.228) -- (7.998,0.228) -- (8.012,0.228) -- (8.026,0.228) -- (8.040,0.228) -- (8.054,0.227) -- (8.068,0.227) -- (8.082,0.227) -- (8.096,0.227) -- (8.110,0.227) -- (8.124,0.227) -- (8.138,0.227) -- (8.152,0.227) -- (8.166,0.226) -- (8.180,0.226) -- (8.194,0.226) -- (8.208,0.226) -- (8.222,0.226) -- (8.236,0.226) -- (8.250,0.225) -- (8.264,0.225) -- (8.278,0.225) -- (8.292,0.225) -- (8.306,0.225) -- (8.320,0.225) -- (8.334,0.224) -- (8.348,0.224) -- (8.362,0.224) -- (8.376,0.224) -- (8.390,0.223) -- (8.404,0.223) -- (8.418,0.223) -- (8.432,0.223) -- (8.446,0.222) -- (8.460,0.222) -- (8.474,0.222) -- (8.488,0.221) -- (8.502,0.221) -- (8.516,0.221) -- (8.530,0.220) -- (8.544,0.220) -- (8.558,0.220) -- (8.572,0.219) -- (8.586,0.219) -- (8.600,0.218) -- (8.614,0.218) -- (8.628,0.218) -- (8.642,0.217) -- (8.656,0.217) -- (8.670,0.216) -- (8.684,0.216) -- (8.698,0.215) -- (8.712,0.215) -- (8.726,0.214) -- (8.740,0.213) -- (8.754,0.213) -- (8.768,0.212) -- (8.782,0.212) -- (8.796,0.211) -- (8.810,0.210) -- (8.824,0.209) -- (8.838,0.209) -- (8.852,0.208) -- (8.866,0.207) -- (8.880,0.206) -- (8.894,0.206) -- (8.908,0.205) -- (8.922,0.204) -- (8.936,0.203) -- (8.950,0.202) -- (8.964,0.201) -- (8.978,0.200) -- (8.992,0.199) -- (9.006,0.198) -- (9.020,0.197) -- (9.034,0.195) -- (9.048,0.194) -- (9.062,0.193) -- (9.076,0.192) -- (9.090,0.190) -- (9.104,0.189) -- (9.118,0.188) -- (9.132,0.186) -- (9.146,0.185) -- (9.160,0.183) -- (9.174,0.181) -- (9.188,0.180) -- (9.202,0.178) -- (9.216,0.176) -- (9.230,0.174) -- (9.244,0.173) -- (9.258,0.171) -- (9.272,0.169) -- (9.286,0.167) -- (9.300,0.164) -- (9.314,0.162) -- (9.328,0.160) -- (9.342,0.158) -- (9.356,0.155) -- (9.370,0.153) -- (9.384,0.150) -- (9.398,0.148) -- (9.412,0.145) -- (9.426,0.142) -- (9.440,0.139) -- (9.454,0.136) -- (9.468,0.133) -- (9.482,0.130) -- (9.496,0.127) -- (9.510,0.123) -- (9.524,0.120) -- (9.538,0.116) -- (9.552,0.113) -- (9.566,0.109) -- (9.580,0.105) -- (9.594,0.101) -- (9.608,0.097) -- (9.622,0.093) -- (9.636,0.088) -- (9.650,0.084) -- (9.664,0.079) -- (9.678,0.075) -- (9.692,0.070) -- (9.706,0.065) -- (9.720,0.060) -- (9.734,0.054) -- (9.748,0.049) -- (9.762,0.044) -- (9.776,0.038) -- (9.790,0.032) -- (9.804,0.026) -- (9.818,0.020) -- (9.832,0.014) -- (9.846,0.007) -- (9.860,0.000) -- (9.874,-0.006) -- (9.888,-0.013) -- (9.902,-0.021) -- (9.916,-0.028) -- (9.930,-0.036) -- (9.944,-0.043) -- (9.958,-0.051) -- (9.972,-0.059) -- (9.986,-0.068) -- (10.000,-0.076) -- (10.014,-0.085) -- (10.028,-0.094) -- (10.042,-0.103) -- (10.056,-0.112) -- (10.070,-0.122) -- (10.084,-0.132) -- (10.098,-0.142) -- (10.112,-0.152) -- (10.126,-0.162) -- (10.140,-0.173) -- (10.154,-0.184) -- (10.168,-0.195) -- (10.182,-0.207) -- (10.196,-0.218) -- (10.210,-0.230) -- (10.224,-0.242) -- (10.238,-0.255) -- (10.252,-0.267) -- (10.266,-0.280) -- (10.280,-0.293) -- (10.294,-0.307) -- (10.308,-0.321) -- (10.322,-0.335) -- (10.336,-0.349) -- (10.350,-0.363) -- (10.364,-0.378) -- (10.378,-0.393) -- (10.392,-0.408) -- (10.406,-0.424) -- (10.420,-0.440) -- (10.434,-0.456) -- (10.448,-0.472) -- (10.462,-0.489) -- (10.476,-0.506) -- (10.490,-0.523) -- (10.504,-0.540) -- (10.518,-0.558) -- (10.532,-0.576) -- (10.546,-0.595) -- (10.560,-0.613) -- (10.574,-0.632) -- (10.588,-0.651) -- (10.602,-0.670) -- (10.616,-0.690) -- (10.630,-0.710) -- (10.644,-0.730) -- (10.658,-0.751) -- (10.672,-0.771) -- (10.686,-0.792) -- (10.700,-0.814) -- (10.714,-0.835) -- (10.728,-0.857) -- (10.742,-0.879) -- (10.756,-0.901) -- (10.770,-0.923) -- (10.784,-0.946) -- (10.798,-0.969) -- (10.812,-0.992) -- (10.826,-1.015) -- (10.840,-1.039) -- (10.854,-1.062) -- (10.868,-1.086) -- (10.882,-1.110) -- (10.896,-1.135) -- (10.910,-1.159) -- (10.924,-1.184) -- (10.938,-1.209) -- (10.952,-1.234) -- (10.966,-1.259) -- (10.980,-1.284) -- (10.994,-1.309) -- (11.008,-1.335) -- (11.022,-1.360) -- (11.036,-1.386) -- (11.050,-1.412) -- (11.064,-1.438) -- (11.078,-1.464) -- (11.092,-1.490) -- (11.106,-1.516) -- (11.120,-1.542) -- (11.134,-1.568) -- (11.148,-1.594) -- (11.162,-1.621) -- (11.176,-1.647) -- (11.190,-1.673) -- (11.204,-1.699) -- (11.218,-1.725) -- (11.232,-1.751) -- (11.246,-1.778) -- (11.260,-1.803) -- (11.274,-1.829) -- (11.288,-1.855) -- (11.302,-1.881) -- (11.316,-1.906) -- (11.330,-1.932) -- (11.344,-1.957) -- (11.358,-1.982) -- (11.372,-2.007) -- (11.386,-2.031) -- (11.400,-2.056) -- (11.414,-2.080) -- (11.428,-2.104) -- (11.442,-2.127) -- (11.456,-2.151) -- (11.470,-2.174) -- (11.484,-2.196) -- (11.498,-2.219) -- (11.512,-2.241) -- (11.526,-2.262) -- (11.540,-2.284) -- (11.554,-2.304) -- (11.568,-2.325) -- (11.582,-2.345) -- (11.596,-2.364) -- (11.610,-2.383) -- (11.624,-2.401) -- (11.638,-2.419) -- (11.652,-2.436) -- (11.666,-2.453) -- (11.680,-2.469) -- (11.694,-2.484) -- (11.708,-2.499) -- (11.722,-2.513) -- (11.736,-2.526) -- (11.750,-2.538) -- (11.764,-2.550) -- (11.778,-2.561) -- (11.792,-2.571) -- (11.806,-2.580) -- (11.820,-2.588) -- (11.834,-2.596) -- (11.848,-2.602) -- (11.862,-2.607) -- (11.876,-2.611) -- (11.890,-2.615) -- (11.904,-2.617) -- (11.918,-2.618) -- (11.932,-2.617) -- (11.946,-2.616) -- (11.960,-2.613) -- (11.974,-2.609) -- (11.988,-2.604) -- (12.002,-2.597) -- (12.016,-2.588) -- (12.030,-2.578) -- (12.044,-2.567) -- (12.058,-2.554) -- (12.072,-2.539) -- (12.086,-2.523) -- (12.100,-2.505) -- (12.114,-2.485) -- (12.128,-2.463) -- (12.142,-2.439) -- (12.156,-2.413) -- (12.170,-2.385) -- (12.184,-2.355) -- (12.198,-2.322) -- (12.212,-2.287) -- (12.226,-2.250) -- (12.240,-2.210) -- (12.254,-2.168) -- (12.268,-2.123) -- (12.282,-2.075) -- (12.296,-2.025) -- (12.310,-1.971) -- (12.324,-1.915) -- (12.338,-1.855) -- (12.352,-1.792) -- (12.366,-1.726) -- (12.380,-1.656) -- (12.394,-1.582) -- (12.408,-1.505) -- (12.422,-1.424) -- (12.436,-1.339) -- (12.450,-1.249) -- (12.464,-1.156) -- (12.478,-1.058) -- (12.492,-0.955) -- (12.506,-0.848) -- (12.520,-0.735) -- (12.534,-0.618) -- (12.548,-0.495) -- (12.562,-0.367) -- (12.576,-0.233) -- (12.590,-0.093) -- (12.604,0.053) -- (12.618,0.205) -- (12.632,0.364) -- (12.646,0.530) -- (12.660,0.702) -- (12.674,0.882) -- (12.688,1.070) -- (12.702,1.265) -- (12.716,1.469) -- (12.730,1.681) -- (12.744,1.901) -- (12.758,2.131) -- (12.772,2.370) -- (12.786,2.618) -- (12.800,2.877);}
\newcommand{\figModelCoordinates}{\coordinate (fB) at (0.467,-0.141); \coordinate (fC) at (5.200,2.581); \coordinate (zA) at (8.775,0.212); \coordinate (zB) at (11.693,-2.483);}
\newcommand{\figModelZoomTickThree}{8.075}
\newcommand{\figModelZoomTickTwo}{9.825}
\newcommand{\figModelZoomTickOne}{11.575}

\begin{tikzpicture}[>=Stealth,font=\small]
  \figModelCoordinates
  % (a) linear scale in rho
  \draw[->,gray] (0,0)--(5.8,0) node[right] {$\rho$};
  \draw[->,gray] (0,-0.45)--(0,3.35);
  \figModelFull{thick}
  \draw[gray] (fB |- 0,0.06)--(fB |- 0,-0.06);
  \node[below] at (fB |- 0,-0.2) {$\rho_{\rm b}$};
  \draw[gray] (fC |- 0,0.06)--(fC |- 0,-0.06) node[below] {$\rho_{\rm c}$};
  \draw[gray,dashed] (fC)--(fC |- 0,0);
  \fill (fB) circle (1.5pt) (fC) circle (1.5pt);
  \node[left] at (fC) {$+$};
  \node at (2.9,-0.8) {(a)};
  % (b) logarithmic scale in rho, near rho_a and rho_b
  \begin{scope}
    \clip (7.2,-2.85) rectangle (12.9,1.05);
    \figModelZoom{thick}
  \end{scope}
  \draw[->,gray] (7.2,0)--(13.0,0) node[right] {$\rho$};
  \draw[->,gray] (7.2,-2.85)--(7.2,1.2);
  \foreach \x/\l in {\figModelZoomTickThree/{10^{-3}},\figModelZoomTickTwo/{10^{-2}},
                     \figModelZoomTickOne/{10^{-1}}}
    {\draw[gray] (\x,0.06)--(\x,-0.06) node[below,font=\scriptsize] {$\l$};}
  \draw[gray,dashed] (zA)--(zA |- 0,0);
  \node[above,font=\scriptsize] at (zA |- 0,0.35) {$\rho_{\rm a}$};
  \draw[gray] (zB |- 0,0.06)--(zB |- 0,-0.06);
  \node[above,font=\scriptsize] at (zB |- 0,0.1) {$\rho_{\rm b}$};
  \fill (zA) circle (1.5pt) (zB) circle (1.5pt);
  \node[above right] at (zA) {$+$};
  \node[below] at (zB) {$-$};
  \node at (10.05,-3.3) {(b)};
\end{tikzpicture}
\caption{The model profile \(D_*+g_{\lambda,\Lambda}(\rho)\) of
\eqref{eq:sharpness-prepared-displacement} for \(C_0=65\), computed from
its formula with \(\lambda=\Lambda_0=\e^{-\pi/2}\) and
\(\Lambda=\Lambda_{\rm b}=\e^{\pi}\). The radii
\(\rho_{\rm a}<\rho_{\rm b}<\rho_{\rm c}\) and the constant \(D_*\) are
chosen as in the proof. (a) Linear scale in \(\rho\). (b) Logarithmic scale
in \(\rho\) near \(\rho_{\rm a}\) and \(\rho_{\rm b}\), with the vertical scale
enlarged. The two model zeros on the compact exit interval lie between
the three marked radii. The fourth sign at a larger central height in
\eqref{eq:sharpness-four-source-signs} gives the third displacement zero.}
\label{fig:sharpness-model}
\end{figure}


\emph{Realizing the additive constant.}
The choices above fix the desired normalized constant \(D_*\).
To realize it, we apply the intermediate value theorem to the
unscaled splitting. By
\eqref{eq:sharpness-endpoint-limit},
\[
 \mathcal S_s(A)\longrightarrow \mathcal S_0(A)
\]
uniformly on a fixed small \(A\)-interval. By
\cref{lem:organizing-splitting-transversality}, choose its endpoints
\(A_{\rm lo}<1<A_{\rm hi}\) so that \(\mathcal S_0\) has opposite
signs there and its only zero in the interval is \(1\).
The continuous function
\[
 \Phi_s(A)=\mathcal S_s(A)-\tfrac12s^2a_+^e(A,s)D_*
\]
converges uniformly to \(\mathcal S_0\). The intermediate value theorem
therefore gives \(A_s\in(A_{\rm lo},A_{\rm hi})\) with
\begin{equation}
\label{eq:sharpness-tuned-splitting}
 \mathcal S_s(A_s)=\tfrac12s^2a_+^e(A_s,s)D_*.
\end{equation}
Every accumulation point of any such choices is a zero of
\(\mathcal S_0\), so \(A_s\to1\). Put \(p_s=p_s(A_s)\).

\emph{Three isolated zeros.}
Let \(h_s(\rho)\) be the common central height of
\eqref{eq:sharpness-common-source} at \(p_s\).
By \eqref{eq:sharpness-prepared-displacement},
\eqref{eq:sharpness-three-model-signs}, and
\eqref{eq:sharpness-tuned-splitting}, with the remainder made
smaller than the fixed margins, \(\Delta_{p_s}\) is positive at
\(h_s(\rho_{\rm a})\) and \(h_s(\rho_{\rm c})\), and negative at
\(h_s(\rho_{\rm b})\). On the fixed radius interval, \(h_s\) is
strictly decreasing by the radial construction above.
Moreover \(\mathcal A_+(\rho)\ge c\log(1/s)-C\), so these three
central heights tend to zero.

The three test radii give two zeros. For a third zero, we need a
negative displacement at a larger central height. We obtain it from
the organizing contraction at a fixed height.
At the organizing parameter, \cref{prop:signed-contraction-new}
gives \(\Delta_0'(h)<0\), while
\(\Delta_0(h)\to0\) as \(h\downarrow0\). Choose a fixed small
\(\bar h>0\) whose endpoints lie in \(\operatorname{int}W_\pm\). Then
\(\Delta_0<0\) on \((0,\bar h]\). For each fixed
\(0<\widetilde h<\bar h\), regular dependence on the parameters at
positive central heights (\cref{lem:regular-transitions}, applied to the
field extended by \(\dot p=0\)) gives
\(\Delta_{p_s}\to\Delta_0\) uniformly on \([\widetilde h,\bar h]\).
Hence \(\Delta_{p_s}<0\) on \([\widetilde h,\bar h]\) for small \(s\);
in particular \(\Delta_{p_s}(\bar h)<0\).

All four central heights belong to \(\cH_{p_s}\), by
\cref{prop:central-displacement-interval-new}, and satisfy
\begin{equation}
\label{eq:sharpness-four-source-signs}
 h_s(\rho_{\rm c})<h_s(\rho_{\rm b})<h_s(\rho_{\rm a})<\bar h,
 \qquad
 \operatorname{sgn}\Delta_{p_s}=(+,-,+,-).
\end{equation}
Since the admissible heights form an interval, there is a zero in
each of the three intervening open intervals.

By \cref{prop:central-displacement-interval-new},
\(\Delta_{p_s}\) is real analytic on
\(\operatorname{int}\cH_{p_s}\), and the four signs in
\eqref{eq:sharpness-four-source-signs} exclude an identically zero
displacement. The three selected zeros are therefore isolated and,
by the correspondence in that proposition, give three distinct
isolated periodic orbits.

Let their central heights be \(h_{j,s}\), \(j=1,2,3\).
The first two are trapped between test heights tending to zero.
For the third, negativity of \(\Delta_{p_s}\) on every fixed
\([\widetilde h,\bar h]\) for small \(s\) excludes a positive
limiting height. Thus all three heights tend to zero, including the
zero between \(h_s(\rho_{\rm a})\) and the fixed height \(\bar h\).

\emph{Convergence of the orbit segments.}
To pass from vanishing central heights to convergence of the orbits,
we compare their graphs near infinity and along the two organizing
branches. Regular flow dependence then treats the finite connecting
arc.
The half-orbits of the smallest test central height
\(h_{{\rm low},s}=h_s(\rho_{\rm c})\) have radius \(O(s)\) in the
thin passage, energy \(O(s^2)\) at \(t_s\), and thereafter lie in
the common cone. Hence, for each fixed small \(\zeta>0\), their
graphs satisfy
\[
 0\le U_\sigma(z;h_{{\rm low},s})\le C\zeta
       \quad(h_{{\rm low},s}\le z\le\zeta),
\]
and converge uniformly to \(U_\sigma^0(z;1)\) on
\([\zeta,z_*]\), by \eqref{eq:sharpness-tail-comparison}.
Every selected zero lies above \(h_{{\rm low},s}\). The exact graph
order \eqref{eq:central-graph-order-new} therefore bounds its energy
above by the graph for \(h_{{\rm low},s}\). Its portion with
\(z\le\zeta\) lies in a neighborhood shrinking to the point at
infinity as \(\zeta\to0\).

To obtain convergence along both organizing branches away from
infinity, fix \(0<\varepsilon<\zeta\). Once
\(h_{j,s}<\varepsilon\), the same order gives, for
\(\zeta\le z\le z_*\),
\[
 U_\sigma(z;\varepsilon,p_s)
 \le U_\sigma(z;h_{j,s},p_s)
 \le U_\sigma(z;h_{{\rm low},s},p_s).
\]
By regular dependence, the left graph converges to the
organizing-parameter graph with central height \(\varepsilon\).
As \(\varepsilon\to0\), that graph converges uniformly on
\([\zeta,z_*]\) to \(U_\sigma^0(z;1)\): at
\(z=K\varepsilon\), for fixed large \(K\), both graphs lie in the
common cone with energy difference \(O(\varepsilon)\), and the
integrable variation coefficient \(O(z^{-1/2})\) propagates this
difference to every fixed height. The two bounds give uniform
convergence of both branches of each periodic orbit away from the
point at infinity.

In particular, the endpoints converge to the organizing endpoints.
The exterior passage lies in finitely many regular flow boxes, so
endpoint convergence and regular dependence of the flow give uniform
convergence of the entire finite connecting arc in the compactified
phase space.

\emph{Hausdorff convergence to the whole graphic.}
We now express these segment estimates in the fixed metric on
\(\mathbb S^2\). Let \(\mathscr P_p\) denote the compactification
map from the parameter-dependent \((\eta,z)\)-chart to \(\mathbb S^2\),
and let \(q_\infty=\mathscr P_{p_0}(0,0)\) be the nilpotent point of the
graphic. Uniform smoothness of these charts gives, on a smaller chart, a
constant \(C_{\mathbb S}\) such that
\begin{equation}
\label{eq:sharpness-sphere-tubular-estimate}
 d_{\mathbb S}\bigl(\mathscr P_p(\eta,z),q_\infty\bigr)
 \le C_{\mathbb S}\bigl(|\eta|+|z|+|p-p_0|\bigr).
\end{equation}
Because \(\eta_\sigma=\sigma\sqrt{2U_\sigma}\), the portions of the
three orbits with \(0\le z\le\zeta\) are at distance at most
\(C_{\mathbb S}'\sqrt{\zeta}+o_s(1)\) from \(q_\infty\), uniformly in
\(s\) after the fixed \(\zeta\) is chosen. The same estimate applies to
the corresponding part of \(\GammaI\) (with \(p=p_0\)).

For the reverse approximation near infinity, the central section point
\(\mathscr P_{p_s}(0,h_{j,s})\) belongs to each periodic orbit and tends to
\(q_\infty\). Thus every point of \(\GammaI\) with
\(0\le z\le\zeta\) is within
\(C_{\mathbb S}'\sqrt{\zeta}+o_s(1)\) of that orbit.
On the remaining portions, the graph convergence on \([\zeta,z_*]\)
and the regular-flow-box convergence of the finite connecting arcs
give both inclusions in the Hausdorff estimate. Therefore, for every
fixed \(\zeta>0\),
\[
 \limsup_{s\downarrow0}
 \distH(\overline{\gamma_{j,s}},\GammaI)
 \le C_{\mathbb S}'\sqrt{\zeta},\qquad j=1,2,3.
\]
Letting \(\zeta\downarrow0\) proves
\[
 \distH(\overline{\gamma_{j,s}},\GammaI)\longrightarrow0,
 \qquad j=1,2,3.
\]
Taking any sequence \(s_n\downarrow0\) completes the proof.
\end{proof}

\begin{proof}[Proof of \cref{cor:exact-local-cyclicity-three}]
By \cref{thm:direct-finite-cyclicity}, \(\Cycl(\GammaI)\le3\).
Let \(V\) be a neighborhood of \((1,C_0,0,0,0)\) and let
\(\varepsilon>0\). For large \(n\), the parameter \(p_n\) of
\cref{prop:three-cycle-lower-bound} lies in \(V\), and its three
limit cycles \(\gamma_{j,n}\) satisfy
\(\distH(\overline{\gamma_{j,n}},\GammaI)<\varepsilon\). Hence
\(N(V,\varepsilon)\ge3\) in \eqref{eq:cyclicity-definition}, and
\(\Cycl(\GammaI)\ge3\).
\end{proof}

\subsection{Hyperbolicity and separation of central heights}

\begin{corollary}[Three hyperbolic limit cycles]
\label{cor:hyperbolic-triple}
For sufficiently small positive \(s\), the three periodic
orbits constructed above are hyperbolic. In increasing order
of their central heights, they are attracting, repelling, and
attracting. For each such \(s\), all three persist, with these
stabilities, on an open neighborhood of \(p_s\) in the original
five-parameter space. Such open sets occur arbitrarily close
to the organizing parameter, with all three orbits arbitrarily
close to the whole graphic.
\end{corollary}

\begin{proof}
We obtain simplicity from the zero bound for the logarithmic slope
ratio: the two gaps between the three displacement zeros already
require two slope zeros. By \cref{prop:direct-buffered-negative},
which applies to \(p_s\) by \eqref{eq:sharpness-coefficient-path},
\(\mathcal L\) has at most two distinct zeros on all of \(\cH_{p_s}\).
By \eqref{eq:direct-completed-log-ratio}, where \(I_{+,h}>0\),
\(\Delta_{p_s,h}\) has the sign of \(\e^{\mathcal L}-1\).
Write the three displacement zeros as \(h_1<h_2<h_3\).
Rolle's theorem gives a zero of \(\mathcal L\) in each of
\((h_1,h_2)\) and \((h_2,h_3)\). If
\(\Delta_{p_s,h}(h_i)=0\) for any \(i\), then \(h_i\) would be a
third distinct zero of \(\mathcal L\). Thus all three zeros are
simple. The same argument excludes a fourth displacement zero.

The signs in \eqref{eq:sharpness-four-source-signs} now imply
\[
 \Delta_{p_s,h}(h_1)<0,\qquad
 \Delta_{p_s,h}(h_2)>0,\qquad
 \Delta_{p_s,h}(h_3)<0.
\]
At each fixed point, differentiation of
\(\cP_{p_s}=I_+^{-1}\circ\Pi_{p_s}\circ I_-\) gives
\begin{equation}
\label{eq:sharpness-return-multipliers}
 \cP_{p_s}'(h_i)=\e^{\mathcal L(h_i)}.
\end{equation}
The two negative signs give multipliers in \((0,1)\), and the
positive sign gives a multiplier greater than one. This proves
hyperbolicity and the stated stabilities.

For fixed \(s>0\), regular dependence
(\cref{lem:regular-transitions}, applied to the field extended by
\(\dot p=0\)) makes these finite transverse returns smooth in all
five parameters, with central height as the section coordinate.
Apply the implicit function theorem near the three simple fixed
points and intersect the resulting parameter neighborhoods.
This gives an open set \(V_s\ni p_s\) on which the three distinct
hyperbolic orbits persist. Their stability is unchanged after
reducing \(V_s\).

Given fixed small parameter and whole-graphic neighborhoods,
first take \(s\) small by the lower-bound construction, then
reduce \(V_s\) to retain the geometric margins.
Within the neighborhood and the Hausdorff condition
of \cref{thm:direct-finite-cyclicity}, these are all the
periodic orbits counted there.
\end{proof}

\begin{corollary}[Separation of the third central height]
\label{cor:third-height-separation}
For the parameters \(p_s\) constructed in the proof of
\cref{prop:three-cycle-lower-bound}, write the three central
heights as \(h_{1,s}<h_{2,s}<h_{3,s}\). Then
\begin{equation}
\label{eq:third-height-separation}
 \lim_{s\downarrow0}s\log\frac{h_{3,s}}{h_s(\rho_{\rm a})}=+\infty.
\end{equation}
The same limit holds with \(h_s(\rho_{\rm a})\) replaced by
\(h_{1,s}\) or \(h_{2,s}\).
\end{corollary}

\begin{proof}
We first place the third zero outside each fixed compact exit interval.
Fix \(0<\rho_{\rm lo}<\rho_{\rm a}\). The central-height construction
\eqref{eq:sharpness-common-source} extends to
\([\rho_{\rm lo},\rho_{\rm c}]\); it is strictly decreasing there and
\(h_s(\rho_{\rm lo})\to0\). Hence, for sufficiently small \(s\),
\[
 h_s(\rho_{\rm c})<h_s(\rho_{\rm lo})<\bar h.
\]
Both outer central heights belong to \(\cH_{p_s}\), as shown in the
proof of \cref{prop:three-cycle-lower-bound}. Since this is an interval,
it contains the image of \([\rho_{\rm lo},\rho_{\rm c}]\) under \(h_s\).

By \eqref{eq:sharpness-coefficient-path}, the compact-primary estimate
in \cref{sec:direct-buffered-compact-primary} applies to the chosen
\(p_s\) on \([\rho_{\rm lo},\rho_{\rm c}]\), after reducing \(s\)
with \(\rho_{\rm lo}\) fixed. It gives at most two zeros on the
corresponding height interval. That interval contains \(h_{1,s}\)
and \(h_{2,s}\), so
\[
 h_{3,s}>h_s(\rho_{\rm lo}).
\]
The count includes endpoints; equivalently, one can first apply it
on a slightly larger fixed compact interval.

We now quantify this separation. Subtracting the equations for the
two central heights and using
\eqref{eq:direct-buffered-primary-action} gives
\[
 s\log\frac{h_s(\rho_{\rm lo})}{h_s(\rho_{\rm a})}
  =L_+(\rho_{\rm a})-L_+(\rho_{\rm lo})+O_{\rho_{\rm lo}}(s).
\]
The common action constant cancels, and the terms \(sM_+^0\)
contribute \(O_{\rho_{\rm lo}}(s)\). Given \(K>0\), first fix
\(\rho_{\rm lo}\) so that
\(L_+(\rho_{\rm a})-L_+(\rho_{\rm lo})>K+1\), and then take \(s\)
small. The preceding inequalities give
\(s\log(h_{3,s}/h_s(\rho_{\rm a}))>K\), proving
\eqref{eq:third-height-separation}. Finally,
\(h_{1,s},h_{2,s}<h_s(\rho_{\rm a})\), which proves the last assertion.
\end{proof}

\begin{remark}[A numerical example]
\label{rem:numerical-three-cycles}
Take \(C_0=65\), so that \(\nu=1\) and \(a=1/2\) in
\eqref{eq:sharpness-coefficient-path}, and take \(s=0.05\) and \(A=1\),
with \(\mu\) given by \eqref{eq:direct-coefficient-bridge}. We computed the
return map to the central section numerically, by an adaptive Runge--Kutta
integration of the field in the chart \eqref{eq:projective-coordinates-new},
with \(\log z\) as height variable, and of \eqref{eq:quadratic-family}
between two crossings of a fixed section, with relative tolerance
\(10^{-11}\).

The accompanying
\href{https://github.com/h-lu/I21-reproof/raw/ccbc54f203ef41d891f4e2c890b595955c58b244/supplementary/three-cycles.zip}{numerical supplement}
contains a standalone Python script, its fixed parameters and tolerances,
execution instructions, and the reference table. It uses a Dormand--Prince
5(4) integrator, logarithmic-height bisection, and the same affine field
between section crossings. For central heights with
\(-120\le\log h\le-3.5\), the computation finds three fixed points:
\begin{center}
% Generated by scripts/numerics_three_cycles.py; do not edit.
\begin{tabular}{@{}ccc@{}}
\hline
\(j\) & \(\log h_j\) & multiplier\\
\hline
\(1\) & \(-93.45\) & \(0.095\)\\
\(2\) & \(-34.94\) & \(9.25\)\\
\(3\) & \(-11.39\) & \(0.289\)\\
\hline
\end{tabular}

\end{center}
Here the multiplier is the derivative of the return map at the fixed point.
In increasing central height the orbits are attracting, repelling and
attracting, as in \cref{cor:hyperbolic-triple}, and the largest central
height is much larger than the other two; compare
\cref{cor:third-height-separation}. The computation illustrates the
construction at these moderate parameter values.
\end{remark}
